%%%%%%%%%%%%%%%%%%%%%%%%%%%%%%%%%%%%%%%%%
% CV of Yunseong Moon -- kept in sync with https://yunseong0518.github.io
% Class: cv.cls (Medium Length Professional CV, LaTeXTemplates.com)
%
% Compile with XeLaTeX (or Tectonic):  xelatex cv.tex   ->  cv.pdf
% UTF-8 encoded.
%%%%%%%%%%%%%%%%%%%%%%%%%%%%%%%%%%%%%%%%%

\documentclass{cv}

\usepackage[left=0.75in,top=0.6in,right=0.75in,bottom=0.6in]{geometry}
\usepackage{fontspec}
\usepackage[hidelinks]{hyperref}
\usepackage{enumitem}
\setlist[itemize]{leftmargin=*, nosep}

% Korean name: use Malgun Gothic on Windows, Noto Sans CJK KR elsewhere (e.g. Overleaf).
\IfFontExistsTF{Malgun Gothic}
  {\newfontfamily\hangulfont{Malgun Gothic}}
  {\newfontfamily\hangulfont{Noto Sans CJK KR}}

\newcommand{\me}{\textbf{Yunseong Moon}}
% Keep hyphenated given names on one line (no "Seung-/Hwan" breaks).
\newcommand{\nm}[1]{\mbox{#1}}

% Numbered reference lists: [J1], [C1], [P1] ...
\newlist{refs}{enumerate}{1}
% Ragged-right so long author lists wrap at name boundaries instead of
% overflowing or splitting a name.
\setlist[refs]{leftmargin=2.7em, labelsep=0.5em, itemsep=4pt, topsep=2pt, before=\raggedright}

\hypersetup{pdftitle={CV - Yunseong Moon}, pdfauthor={Yunseong Moon}}

\name{Yunseong Moon ({\hangulfont 문윤성})}

% cv.cls prints the address lines in the order 3rd, 1st, 2nd call, so the line
% meant to appear first (affiliation) is given last.
\address{
\href{mailto:yunseong0518@postech.ac.kr}{yunseong0518@postech.ac.kr} \\
\href{https://yunseong0518.github.io}{yunseong0518.github.io}
}
\address{
\href{https://scholar.google.com/citations?user=YwKp8dIAAAAJ}{Google Scholar} \\
\href{https://github.com/yunseong0518}{github.com/yunseong0518} \\
\href{https://www.linkedin.com/in/yunseongmoon}{linkedin.com/in/yunseongmoon}
}
\address{
Ph.D. Student, Computer Science and Engineering, POSTECH, Republic of Korea
}

\begin{document}

%----------------------------------------------------------------------------------------
\begin{rSection}{Research Interests}
\begin{itemize}
\item Computational imaging: polarimetric and hyperspectral imaging systems
\item Appearance acquisition and modeling (polarimetric/spectral BRDFs)
\item Inverse rendering and 3D imaging
\end{itemize}
\end{rSection}

%----------------------------------------------------------------------------------------
\begin{rSection}{Education}

{\bf Pohang University of Science and Technology (POSTECH)} \hfill {\em Feb 2024 -- Present}\\
M.S.--Ph.D. Integrated Program, Computer Science and Engineering\\
Computer Graphics Lab, Advisor: Prof. \href{https://www.shbaek.com}{Seung-Hwan Baek}

{\bf Pohang University of Science and Technology (POSTECH)} \hfill {\em Feb 2019 -- Feb 2024}\\
B.S., Computer Science and Engineering \hfill {\em Magna Cum Laude}

\end{rSection}

%----------------------------------------------------------------------------------------
\begin{rSection}{Publications}
{\small * Equal contribution}

\textbf{Journal Papers}
\begin{refs}[label={[J\arabic*]}]
\item Seokjun Choi*, \me*, Kaizhang Kang, \nm{Hoon-Gyu Chung}, \nm{Jin-Nyeong Kim}, Giljoo Nam, \nm{Seung-Hwan Baek},
\emph{Snapshot Polarimetric Display Inverse Rendering},
\textbf{ACM Transactions on Graphics} (Proc.\ \textbf{SIGGRAPH Asia} 2026).
\end{refs}

\textbf{Conference Papers}
\begin{refs}[label={[C\arabic*]}]
\item Yonghee Oh, Ryota Maeda, \me, \nm{Seung-Hwan Baek},
\emph{Real-world Polarimetric Environment Map Dataset},
\textbf{SIGGRAPH Asia} 2026.

\item Suhyun Shin, \me, Ryota Maeda, David B. Lindell, Kiriakos N. Kutulakos, \nm{Seung-Hwan Baek},
\emph{Broadband Hyperspectral 3D Imaging under Dispersed Structured Light},
\textbf{SIGGRAPH} 2026.

\item \me, Ryota Maeda, Suhyun Shin, Inseung Hwang, Youngchan Kim, Min H. Kim, \nm{Seung-Hwan Baek},
\emph{Hyperspectral Polarimetric BRDFs of Real-world Materials},
\textbf{SIGGRAPH Asia} 2025.

\item Ryota Maeda, \me, \nm{Seung-Hwan Baek},
\emph{Event Ellipsometer: Event-based Mueller-Matrix Video Imaging},
\textbf{CVPR} 2025, \textbf{Highlight}.

\item Yujin Jeon*, Eunsue Choi*, Youngchan Kim, \me, Khalid Omer, Felix Heide, \nm{Seung-Hwan Baek},
\emph{Spectral and Polarization Vision: Spectro-polarimetric Real-world Dataset},
\textbf{CVPR} 2024, \textbf{Highlight}.
\end{refs}

\textbf{Posters}
\begin{refs}[label={[P\arabic*]}]
\item \nm{Oscar Pueyo-Ciutad}, \nm{Guillermo Enguita-Lahoz}, \me, Ryota Maeda, \nm{Seung-Hwan Baek}, \nm{Albert Redo-Sanchez}, Diego Gutierrez,
\emph{Transient Polarimetry},
\textbf{SIGGRAPH} 2026 Posters.
\end{refs}

\end{rSection}

%----------------------------------------------------------------------------------------
\begin{rSection}{Work Experience}
\begin{itemize}
\item Part-time Software Engineer, Sketchsoft Inc., Seoul, Republic of Korea \hfill {\em Sep 2022 -- Aug 2023}
\item Full-time Software Engineer, Sketchsoft Inc., Seoul, Republic of Korea \hfill {\em Jan 2022 -- Jul 2022}
\item Software Engineer Intern (Full-time), DEMO, Pohang, Republic of Korea \hfill {\em Dec 2020 -- Feb 2021}
\end{itemize}
\end{rSection}

%----------------------------------------------------------------------------------------
\begin{rSection}{Honors and Awards}
\begin{itemize}
\item BK21 Outstanding Paper Award (1st Prize) \hfill {\em Jan 2026}
\end{itemize}
\end{rSection}

%----------------------------------------------------------------------------------------
\begin{rSection}{Teaching Experience}
\begin{itemize}
\item Teaching Assistant, Youth AI$\cdot$Data Science Program -- Computer Vision \hfill {\em Nov -- Dec 2025}
\item Teaching Assistant, CSED520: Computational Imaging, POSTECH \hfill {\em Sep -- Dec 2025}
\item Teaching Assistant, POSCO AI Program -- Computer Vision \hfill {\em Jun 2024}
\end{itemize}
\end{rSection}

%----------------------------------------------------------------------------------------
\begin{rSection}{Professional Service}
\begin{itemize}
\item Reviewer: ACM Transactions on Graphics (TOG), IEEE Transactions on Visualization and Computer Graphics (TVCG), Asian Conference on Computer Vision (ACCV)
\end{itemize}
\end{rSection}

%----------------------------------------------------------------------------------------
\begin{rSection}{Patents}
\begin{itemize}
\item Method and Apparatus for 3D Modeling, Korea Patent No.\ 10-2023-0059316 \hfill {\em Registered May 2023}
\end{itemize}
\end{rSection}

\end{document}
